
\begin{definition}[Retained coordinates for the vertical-sector algebra]
    \label{def:mpdo_retained_vertical_sector_coordinates}
    \lean{MPOTensor.verticalSectorFinEquiv}
    \lean{MPOTensor.verticalSectorBlockDiagonal}
    \lean{MPOTensor.verticalSectorBlocks}
    \lean{MPOTensor.verticalSectorBlockProjection}
    \lean{MPOTensor.normalizedRetainedVerticalSectorEmbedding}
    \lean{MPOTensor.retainedVerticalSectorPartialTrace}
    \leanok
    \uses{def:mpdo_normalized_vertical_sector_maps}
    Retain the coordinates $(\alpha,a,i)$, where $\alpha$ labels a simple
    sector, $a$ labels a diagonal entry of $\mu_\alpha$, and
    $0\leq i<d_\alpha$.  The family $(Y_\alpha)_\alpha$ is represented on
    this space by
    \[
        \bigoplus_\alpha Y_\alpha.
    \]
    Write $\iota$ for this inclusion, $\pi$ for extraction of the diagonal
    summands, and $P=\iota\pi$.  Composing $\iota$ and $\pi$ with $R_\mu$ and
    $\widetilde R_\mu$ defines their retained-coordinate forms.
\end{definition}

\begin{lemma}[Extraction after vertical-sector inclusion]
    \label{lem:mpdo_retained_vertical_sector_extraction}
    \lean{MPOTensor.verticalSectorBlocks_comp_verticalSectorBlockDiagonal}
    \leanok
    \uses{def:mpdo_retained_vertical_sector_coordinates}
    The diagonal-summand extraction is a left inverse of the inclusion:
    \[
        \pi\iota=\id.
    \]
\end{lemma}

\begin{proof}\leanok
    Restricting $\bigoplus_\beta Y_\beta$ to the rows and columns with sector
    label $\alpha$ gives $Y_\alpha$.
\end{proof}

\begin{lemma}[Entries retained by the vertical-sector projection]
    \label{lem:mpdo_retained_vertical_sector_projection_same}
    \lean{MPOTensor.verticalSectorBlockProjection_apply_same}
    \leanok
    \uses{def:mpdo_retained_vertical_sector_coordinates}
    For every retained-coordinate matrix $Z$,
    \[
        P(Z)_{(\alpha,a,i),(\alpha,b,j)}
          =Z_{(\alpha,a,i),(\alpha,b,j)}.
    \]
\end{lemma}

\begin{proof}\leanok
    Both indices lie in the same extracted diagonal summand.
\end{proof}

\begin{lemma}[Entries removed by the vertical-sector projection]
    \label{lem:mpdo_retained_vertical_sector_projection_distinct}
    \lean{MPOTensor.verticalSectorBlockProjection_apply_ne}
    \leanok
    \uses{def:mpdo_retained_vertical_sector_coordinates}
    If $\alpha\ne\beta$, then
    \[
        P(Z)_{(\alpha,a,i),(\beta,b,j)}=0.
    \]
\end{lemma}

\begin{proof}\leanok
    The two indices belong to distinct diagonal summands.
\end{proof}

\begin{lemma}[Idempotence of the vertical-sector projection]
    \label{lem:mpdo_retained_vertical_sector_projection_idempotent}
    \lean{MPOTensor.verticalSectorBlockProjection_comp_self}
    \leanok
    \uses{lem:mpdo_retained_vertical_sector_extraction}
    The projection onto the vertical-sector diagonal summands is idempotent:
    \[
        P^2=P.
    \]
\end{lemma}

\begin{proof}\leanok
    The identity $\pi\iota=\id$ gives
    $P^2=\iota\pi\iota\pi=\iota\pi$.
\end{proof}

\begin{lemma}[Retraction in retained vertical-sector coordinates]
    \label{lem:mpdo_retained_vertical_sector_retraction}
    \lean{MPOTensor.retainedVerticalSectorPartialTrace_comp_normalizedEmbedding}
    \leanok
    \uses{def:mpdo_retained_vertical_sector_coordinates,
      lem:mpdo_retained_vertical_sector_extraction,
      lem:mpdo_vertical_sector_retraction}
    Suppose that every multiplicity space is nonzero and every diagonal entry
    of $\mu_\alpha$ is positive.  Then the retained-coordinate partial trace
    is a left inverse of the normalized retained-coordinate embedding.
\end{lemma}

\begin{proof}\leanok
    First extract the diagonal summands, then apply
    $\widetilde R_\mu R_\mu=\id$.
\end{proof}

\begin{lemma}[Weighted vertical contraction in retained coordinates]
    \label{lem:mpdo_retained_vertical_weighted_contraction}
    \lean{MPOTensor.contractBondMatrix_verticalAssembledTensor_eq_sectorBlockDiagonal}
    \leanok
    \uses{def:mpdo_retained_vertical_sector_coordinates,
      def:mpdo_vertical_bond_matrix_contraction}
    For every horizontal bond matrix $X$, contraction of
    $\bigoplus_\alpha\mu_\alpha\otimes M_\alpha$ gives
    \[
        \bigoplus_\alpha\mu_\alpha\otimes M_\alpha(X).
    \]
\end{lemma}

\begin{proof}\leanok
    Flattening the index triple $(\alpha,a,i)$ gives the entry formula
    \[
        B(X)_{(\alpha,a,i),(\beta,b,j)}
          =\delta_{\alpha\beta}\delta_{ab}\,
            \mu_{\alpha,a}M_\alpha(X)_{ij}.
    \]
\end{proof}

\begin{definition}[Bond-matrix contraction of a vertical tensor]
    \label{def:mpdo_vertical_bond_matrix_contraction}
    \lean{MPSTensor.contractBondMatrix}
    \leanok
    Let $A$ be a tensor whose physical label is a pair $(a,b)$ of horizontal
    bond indices.  For a bond matrix $X$, define
    \[
        A(X)=\sum_{a,b}X_{ba}A_{ab}.
    \]
    The order of the entries of $X$ is chosen so that
    $A(X)_{ij}=\tr(M^{ij}X)$ when
    $A_{ab,ij}=M^{ij}_{ab}$.
\end{definition}

\begin{lemma}[Contraction under fixed left and right multiplication]
    \label{lem:mpdo_bond_contraction_composition}
    \lean{MPSTensor.contractBondMatrix_mul_mul}
    \leanok
    \uses{def:mpdo_vertical_bond_matrix_contraction}
    Let $L$ and $K$ be fixed matrices of compatible sizes.  Then
    \[
        \left((a,b)\mapsto L A_{ab}K\right)(X)=L A(X)K.
    \]
\end{lemma}

\begin{proof}\leanok
    Distributivity of matrix multiplication over finite sums gives
    \[
        \sum_{a,b}X_{ba}\left(LA_{ab}K\right)
          =L\left(\sum_{a,b}X_{ba}A_{ab}\right)K.
    \]
\end{proof}

\begin{lemma}[Contraction of a weighted block sum]
    \label{lem:mpdo_bond_contraction_block_sum}
    \lean{MPSTensor.contractBondMatrix_toTensorFromBlocks}
    \leanok
    \uses{def:mpdo_vertical_bond_matrix_contraction}
    For a weighted block-diagonal tensor
    $A=\bigoplus_\alpha\mu_\alpha A_\alpha$, one has
    \[
        A(X)=\bigoplus_\alpha\mu_\alpha A_\alpha(X).
    \]
\end{lemma}

\begin{proof}\leanok
    Since each tensor letter is block diagonal,
    \[
        \sum_{a,b}X_{ba}
          \left(\bigoplus_\alpha\mu_\alpha A_\alpha\right)_{ab}
          =\bigoplus_\alpha\mu_\alpha
            \sum_{a,b}X_{ba}(A_\alpha)_{ab}
          =\bigoplus_\alpha\mu_\alpha A_\alpha(X).
    \]
\end{proof}

\begin{lemma}[Vertical contraction and physical closure]
    \label{lem:mpdo_vertical_bond_contraction_phys_close}
    \lean{MPOTensor.contractBondMatrix_verticalTensor_eq_physClose1}
    \leanok
    \uses{def:phys_close, def:mpdo_vertical_bond_matrix_contraction}
    For the vertically viewed tensor
    $\widetilde M_{ab,ij}=M^{ij}_{ab}$, bond-matrix contraction is the
    one-site physical closure:
    \[
        \widetilde M(X)=M(X).
    \]
\end{lemma}

\begin{proof}\leanok
    For every pair of physical indices,
    \[
        \widetilde M(X)_{ij}
          =\sum_{a,b}X_{ba}M^{ij}_{ab}
          =\tr(M^{ij}X)
          =M(X)_{ij}.
    \]
\end{proof}

\begin{lemma}[Contraction of the assembled vertical tensor]
    \label{lem:mpdo_vertical_bond_contraction_assembled}
    \lean{MPOTensor.contractBondMatrix_verticalAssembledTensor}
    \leanok
    \uses{lem:mpdo_bond_contraction_block_sum}
    The contraction of the assembled vertical tensor is the weighted block
    diagonal of the sector contractions:
    \[
        \left(\bigoplus_\alpha\mu_\alpha\otimes M_\alpha\right)(X)
          =\bigoplus_\alpha\mu_\alpha\otimes M_\alpha(X).
    \]
\end{lemma}

\begin{proof}\leanok
    Applying the preceding lemma to the block form
    $\bigoplus_\alpha\mu_\alpha\otimes M_\alpha$ of the assembled tensor gives
    \[
        \left(\bigoplus_\alpha\mu_\alpha\otimes M_\alpha\right)(X)
          =\bigoplus_\alpha\mu_\alpha\otimes M_\alpha(X).
    \]
\end{proof}

\begin{lemma}[Contraction of a letterwise forward identity]
    \label{lem:mpdo_vertical_bond_contraction_forward}
    \lean{MPOTensor.mul_physClose1_mul_conjTranspose_of_vertical_forward}
    \leanok
    \uses{def:phys_close, def:mpdo_vertical_bond_matrix_contraction,
      lem:mpdo_bond_contraction_composition,
      lem:mpdo_vertical_bond_contraction_phys_close}
    Let $B_{ab}$ be an arbitrary vertical tensor.  If
    \[
        U\widetilde M_{ab}U^\dagger=B_{ab}
    \]
    for every pair $(a,b)$, then, for every bond matrix $X$,
    \[
        U M(X)U^\dagger=B(X).
    \]
\end{lemma}

\begin{proof}\leanok
    By the definition of bond-matrix contraction,
    \[
        U M(X)U^\dagger
          =U\left(\sum_{a,b}X_{ba}\widetilde M_{ab}\right)U^\dagger
          =\sum_{a,b}X_{ba}
            \left(U\widetilde M_{ab}U^\dagger\right)
          =\sum_{a,b}X_{ba}B_{ab}
          =B(X).
    \]
\end{proof}

\begin{lemma}[Contraction of a letterwise reconstruction identity]
    \label{lem:mpdo_vertical_bond_contraction_reconstruction}
    \lean{MPOTensor.physClose1_eq_of_vertical_reconstruction}
    \leanok
    \uses{def:phys_close, def:mpdo_vertical_bond_matrix_contraction,
      lem:mpdo_bond_contraction_composition,
      lem:mpdo_vertical_bond_contraction_phys_close}
    Let $B_{ab}$ be an arbitrary vertical tensor.  If
    \[
        \widetilde M_{ab}=U^\dagger B_{ab}U
    \]
    for every pair $(a,b)$, then, for every bond matrix $X$,
    \[
        M(X)=U^\dagger B(X)U.
    \]
\end{lemma}

\begin{proof}\leanok
    Contracting the letterwise identity gives
    \[
        M(X)
          =\sum_{a,b}X_{ba}\widetilde M_{ab}
          =\sum_{a,b}X_{ba}\left(U^\dagger B_{ab}U\right)
          =U^\dagger\left(\sum_{a,b}X_{ba}B_{ab}\right)U
          =U^\dagger B(X)U.
    \]
\end{proof}

\begin{lemma}[Contraction of the one-site vertical canonical form]
    \label{lem:mpdo_vertical_bond_matrix_contraction}
    \lean{MPOTensor.mul_physClose1_mul_conjTranspose_of_verticalAssembledTensor}
    \lean{MPOTensor.physClose1_eq_of_verticalAssembledTensor_reconstruction}
    \leanok
    \uses{def:vertical_cf_mpdo, def:phys_close,
      def:mpdo_vertical_bond_matrix_contraction,
      lem:mpdo_vertical_bond_contraction_forward,
      lem:mpdo_vertical_bond_contraction_reconstruction,
      lem:mpdo_vertical_bond_contraction_assembled}
    Suppose that
    \[
        U\widetilde M_{ab}U^\dagger
          =\left(\bigoplus_\alpha\mu_\alpha\otimes M_\alpha\right)_{ab}
    \]
    is a one-site vertical canonical-form identity.  Contracting against an
    arbitrary bond matrix $X$ gives
    \[
        U M(X)U^\dagger
          =\bigoplus_\alpha\mu_\alpha\otimes M_\alpha(X).
    \]
    If the reverse letterwise identity is also given, then
    \[
        M(X)=U^\dagger
          \left(\bigoplus_\alpha\mu_\alpha\otimes M_\alpha(X)\right)U.
    \]
    These are the contracted canonical-form identities used in
    \cite[Appendix~C.4, lines~1955--1979]{Cirac2016MPDO_arXiv}.
\end{lemma}

\begin{proof}\leanok
    Since $A(X)_{ij}=\sum_{a,b}X_{ba}A_{ab,ij}$ and
    $\widetilde M_{ab,ij}=M^{ij}_{ab}$, one has
    \[
        M(X)=\sum_{a,b}X_{ba}\widetilde M_{ab}.
    \]
    Distributivity of fixed left and right multiplication over finite sums
    gives
    \[
        U\left(\sum_{a,b}X_{ba}\widetilde M_{ab}\right)U^\dagger
          =\sum_{a,b}X_{ba}
            \left(U\widetilde M_{ab}U^\dagger\right).
    \]
    Substituting the forward hypothesis into this equality gives
    \[
        U M(X)U^\dagger
          =\sum_{a,b}X_{ba}
            \left(\bigoplus_\alpha\mu_\alpha\otimes M_\alpha\right)_{ab}
          =\left(\bigoplus_\alpha\mu_\alpha\otimes M_\alpha\right)(X).
    \]
    Applied blockwise, the same calculation gives
    \[
        \left(\bigoplus_\alpha \mu_\alpha\otimes M_\alpha\right)(X)
          =\bigoplus_\alpha \mu_\alpha\otimes M_\alpha(X).
    \]
    If the reverse letterwise identity is given, the same calculation yields
    \[
    \begin{aligned}
        M(X)
          &=\sum_{a,b}X_{ba}\widetilde M_{ab} \\
          &=U^\dagger\left(
              \sum_{a,b}X_{ba}
              \left(\bigoplus_\alpha\mu_\alpha\otimes M_\alpha\right)_{ab}
            \right)U \\
          &=U^\dagger\left(
              \bigoplus_\alpha\mu_\alpha\otimes M_\alpha(X)
            \right)U.
    \end{aligned}
    \]
\end{proof}

\begin{definition}[Four-site physical closure]
    \label{def:mpdo_four_site_physical_closure}
    \lean{finFourArrowEquiv, MPOTensor.physClose4}
    \leanok
    For a virtual matrix $X$, the four-site physical closure of $K$ is the
    operator $K_4(X)$ whose matrix coefficients are
    \[
        K_4(X)_{(i_0,(i_1,(i_2,i_3))),(j_0,(j_1,(j_2,j_3)))}
          =\tr\!\left(
            K^{i_0j_0}K^{i_1j_1}K^{i_2j_2}K^{i_3j_3}X\right).
    \]
\end{definition}

\begin{lemma}[Coordinates of the four-site regrouping]
    \label{lem:fin_four_arrow_equiv_coordinates}
    \lean{finFourArrowEquiv_apply, finFourArrowEquiv_symm_apply}
    \leanok
    The canonical four-coordinate regrouping and its inverse satisfy
    \[
        R_4(\sigma)
          =(\sigma_0,(\sigma_1,(\sigma_2,\sigma_3))),
        \qquad
        R_4^{-1}(i_0,(i_1,(i_2,i_3)))
          =(i_0,i_1,i_2,i_3).
    \]
\end{lemma}

\begin{proof}\leanok
    Both identities follow from the definition of the canonical
    right-associated regrouping.
\end{proof}

\begin{theorem}[General and right-associated four-site closures]
    \label{thm:mpdo_phys_close_four}
    \lean{MPOTensor.physCloseN_four_eq_physClose4}
    \leanok
    \uses{def:phys_close, def:mpdo_four_site_physical_closure}
    Under the canonical right-associated identification of four-site
    configurations with quadruples of physical indices,
    \[
        K_4(X)_{(i_0,(i_1,(i_2,i_3))),(j_0,(j_1,(j_2,j_3)))}
          =\tr\!\left(K^{i_0j_0}K^{i_1j_1}K^{i_2j_2}K^{i_3j_3}X\right).
    \]
\end{theorem}

\begin{proof}\leanok
    Expand the general four-site closure in the right-associated coordinates.
\end{proof}

\begin{theorem}[One-site closure of the two-site blocking]
    \label{thm:mpdo_one_site_closure_two_site_blocking}
    \lean{MPOTensor.physClose1_blockTwo_eq_physClose2}
    \lean{MPOTensor.physClose1_blockTwo_eq_physClose2_apply}
    \lean{MPOTensor.physClose2_eq_physClose1_blockTwo_apply}
    \leanok
    \uses{def:phys_close, def:mpdo_two_site_blocking}
    After identifying one blocked physical index with two original indices,
    \[
        \bigl(K^{[2]}\bigr)_1(X)=K_2(X).
    \]
\end{theorem}

\begin{proof}\leanok
    Expand one blocked tensor matrix and decode its two physical indices.
\end{proof}

\begin{definition}[Physical maps in vertical canonical coordinates]
    \label{def:mpdo_vertical_coordinate_map_transport}
    \lean{MPOTensor.transportTToVerticalCoordinates}
    \lean{MPOTensor.transportSToVerticalCoordinates}
    \leanok
    \uses{def:mpdo_two_site_blocking}
    Let $U_1$ and $U_2$ be the physical coordinate changes in vertical
    canonical forms of $M$ and its two-site blocking, respectively.  Let
    $J$ be the canonical relabelling from matrices indexed by pairs of
    one-site physical indices to matrices indexed by one blocked physical
    index.  For the refinement and coarse-graining maps $T$ and $S$, define

    \[
        T_{\mathrm v}(Z)
          =U_2 J\!\left(T(U_1^\dagger ZU_1)\right)U_2^\dagger,
        \qquad
        S_{\mathrm v}(W)
          =U_1 S\!\left(J^{-1}(U_2^\dagger WU_2)\right)U_1^\dagger.
    \]
    These are the physical maps after absorbing the two vertical coordinate
    changes, as in
    \cite[Appendix~C.4, lines~1955--1979]{Cirac2016MPDO_arXiv}.  They are
    intermediate maps on the retained physical coordinate spaces; the
    normalized sector maps $R_1,R_2,\widetilde R_1,\widetilde R_2$ have not
    yet been composed with them.
\end{definition}

\begin{theorem}[Action of the transported physical maps]
    \label{thm:mpdo_vertical_coordinate_map_transport}
    \lean{MPOTensor.transportTToVerticalCoordinates_contractBondMatrix}
    \lean{MPOTensor.transportSToVerticalCoordinates_contractBondMatrix}
    \leanok
    \uses{def:mpdo_vertical_coordinate_map_transport,
      lem:mpdo_vertical_bond_contraction_forward,
      lem:mpdo_vertical_bond_contraction_reconstruction,
      thm:mpdo_one_site_closure_two_site_blocking}
    For a bond matrix $X$, write $\mathcal M^{(1)}(X)$ and
    $\mathcal M^{(2)}(X)$ for the one-site and two-site physical closures,
    respectively, and write
    \[
        B_1(X)=\bigoplus_\alpha\mu_\alpha\otimes M_\alpha(X),
        \qquad
        B_2(X)=\bigoplus_\beta\nu_\beta\otimes A_\beta(X)
    \]
    for the contracted one-site and two-site vertical canonical forms.  If

    \[
        T(\mathcal M^{(1)}(X))=\mathcal M^{(2)}(X),
        \qquad
        S(\mathcal M^{(2)}(X))=\mathcal M^{(1)}(X),
    \]
    then
    \[
        T_{\mathrm v}(B_1(X))=B_2(X),
        \qquad
        S_{\mathrm v}(B_2(X))=B_1(X).
    \]
    These are the physical-coordinate identities which, after composition
    with $R_1,R_2,\widetilde R_1,\widetilde R_2$, give the two displayed
    identities for $\widetilde T$ and $\widetilde S$ in
    \cite[Appendix~C.4, lines~1955--1979]{Cirac2016MPDO_arXiv}.
\end{theorem}

\begin{proof}\leanok
    The reverse one-site canonical-form identity gives
    \[
        U_1^\dagger B_1(X)U_1=\mathcal M^{(1)}(X).
    \]
    Applying $T$, relabelling its two-site output by $J$, and using the
    forward blocked canonical-form identity gives
    \[
    \begin{aligned}
        T_{\mathrm v}(B_1(X))
          &=U_2J(T(\mathcal M^{(1)}(X)))U_2^\dagger \\
          &=U_2J(\mathcal M^{(2)}(X))U_2^\dagger
           =B_2(X).
    \end{aligned}
    \]
    The reverse blocked canonical-form identity and the equation for $S$
    give, in the other direction,
    \[
    \begin{aligned}
        S_{\mathrm v}(B_2(X))
          &=U_1S(\mathcal M^{(2)}(X))U_1^\dagger \\
          &=U_1\mathcal M^{(1)}(X)U_1^\dagger
           =B_1(X).
    \end{aligned}
    \]
\end{proof}

\begin{definition}[Refinement and coarse-graining on vertical sectors]
    \label{def:mpdo_transported_vertical_sector_maps}
    \lean{MPOTensor.transportedVerticalSectorT}
    \lean{MPOTensor.transportedVerticalSectorS}
    \leanok
    \uses{def:mpdo_normalized_vertical_sector_maps,
      def:mpdo_retained_vertical_sector_coordinates,
      def:mpdo_vertical_coordinate_map_transport}
    Define maps between the one-site and two-site simple-sector algebras by
    \[
        \widetilde T=\widetilde R_2T_{\mathrm v}R_1,
        \qquad
        \widetilde S=\widetilde R_1S_{\mathrm v}R_2.
    \]
    The partial trace and diagonal-summand extraction define these maps on the
    full matrix algebras.  Their agreement with the normalized maps in
    \cite[Appendix~C.4, lines~1972--1979]{Cirac2016MPDO_arXiv} is asserted on
    the weighted block-diagonal image.  At this stage there is no assertion
    that the image is preserved, that the two maps are inverse, completely
    positive, or trace preserving, or that the sector labels have been
    matched.  The final renormalization fixed-point conclusion is separate.
\end{definition}

\begin{definition}[Positive vertical-sector families and total trace]
    \label{def:mpdo_positive_vertical_sector_family}
    \lean{MPOTensor.IsVerticalSectorPosSemidef}
    \lean{MPOTensor.verticalSectorTrace}
    \leanok
    \uses{def:mpdo_retained_vertical_sector_coordinates}
    A vertical-sector family $X=(X_\alpha)_\alpha$ is positive when every
    $X_\alpha$ is positive semidefinite.  Its total sector trace is
    \[
        \Trsec(X)
          =\sum_\alpha\tr(X_\alpha).
    \]
\end{definition}

\begin{lemma}[Positivity and trace in normalized sector coordinates]
    \label{lem:mpdo_normalized_vertical_sector_positive_trace}
    \lean{MPOTensor.verticalMultiplicityTrace_pos}
    \lean{MPOTensor.normalizedVerticalSectorEmbedding_posSemidef}
    \lean{MPOTensor.normalizedRetainedVerticalSectorEmbedding_posSemidef}
    \lean{MPOTensor.verticalSectorPartialTrace_posSemidef}
    \lean{MPOTensor.retainedVerticalSectorPartialTrace_posSemidef}
    \lean{MPOTensor.trace_verticalSectorBlockDiagonal}
    \lean{MPOTensor.trace_verticalSectorBlockProjection}
    \lean{MPOTensor.verticalSectorTrace_retainedVerticalSectorPartialTrace}
    \lean{MPOTensor.trace_normalizedRetainedVerticalSectorEmbedding}
    \leanok
    \uses{def:mpdo_positive_vertical_sector_family,
      def:mpdo_retained_vertical_sector_coordinates}
    Suppose that every multiplicity space is nonzero and every diagonal entry
    of $\mu_\alpha$ is positive.  The normalized embedding $R_\mu$ sends a
    positive sector family to a positive retained matrix and satisfies
    \[
        \tr(R_\mu(X))
          =\Trsec(X).
    \]
    Conversely, the sectorwise partial trace sends a positive retained matrix
    to a positive sector family and satisfies
    \[
        \Trsec(\widetilde R_\mu(Z))
          =\tr(Z).
    \]
    These trace identities follow by summing the traces of the diagonal
    sectors; deleting the off-diagonal sector entries does not change the
    trace.
\end{lemma}

\begin{proof}\leanok
    A positive diagonal weight matrix has positive trace, and its Kronecker
    product with each positive $X_\alpha$ is positive.  Partial trace preserves
    positivity.  The two trace formulas follow from
    $\tr(A\otimes B)=\tr(A)\tr(B)$
    and the normalization by $\tr(\mu_\alpha)$.
\end{proof}

\begin{lemma}[Normalized embedding of multiplicity-trace-scaled sectors]
    \label{lem:mpdo_normalized_retained_vertical_sector_embedding}
    \lean{MPOTensor.normalizedRetainedVerticalSectorEmbedding_trace_smul}
    \leanok
    \uses{def:mpdo_retained_vertical_sector_coordinates}
    Suppose that every multiplicity space is nonzero and every diagonal entry
    of $\mu_\alpha$ is positive.  For
    $m_\alpha=\tr(\mu_\alpha)$,
    \[
        R_\mu\!\left(\bigoplus_\alpha m_\alpha X_\alpha\right)
          =\bigoplus_\alpha\mu_\alpha\otimes X_\alpha.
    \]
\end{lemma}

\begin{proof}\leanok
    In each sector,
    \[
        R_\mu(m_\alpha X_\alpha)
        =m_\alpha^{-1}m_\alpha(\mu_\alpha\otimes X_\alpha)
        =\mu_\alpha\otimes X_\alpha.
    \]
\end{proof}

\begin{lemma}[Partial trace of weighted vertical sectors]
    \label{lem:mpdo_retained_vertical_sector_partial_trace}
    \lean{MPOTensor.retainedVerticalSectorPartialTrace_weightedBlockDiagonal}
    \leanok
    \uses{def:mpdo_retained_vertical_sector_coordinates}
    For $m_\alpha=\tr(\mu_\alpha)$,
    \[
        \widetilde R_\mu\!\left(
          \bigoplus_\alpha\mu_\alpha\otimes X_\alpha\right)
          =\bigoplus_\alpha m_\alpha X_\alpha.
    \]
\end{lemma}

\begin{proof}\leanok
    Taking the partial trace over the multiplicity factor gives
    \[
        \widetilde R_\mu(\mu_\alpha\otimes X_\alpha)
        =\tr(\mu_\alpha)X_\alpha
        =m_\alpha X_\alpha.
    \]
\end{proof}

\begin{lemma}[Trace under coisometric compression]
    \label{lem:mpdo_coisometric_compression_trace}
    \lean{Matrix.isOrthogonalProjection_conjTranspose_mul_of_mul_conjTranspose_eq_one}
    \lean{Matrix.PosSemidef.trace_mul_mul_conjTranspose_le_of_mul_conjTranspose_eq_one}
    \lean{Matrix.PosSemidef.trace_mul_mul_conjTranspose_eq_iff_of_mul_conjTranspose_eq_one}
    \lean{Matrix.trace_conjTranspose_mul_mul_of_mul_conjTranspose_eq_one}
    \leanok
    Let $U:\mathbb C^n\to\mathbb C^r$ satisfy $UU^\dagger=1$.  If $Z\geq0$,
    then
    \[
        \tr(UZU^\dagger)\leq\tr(Z),
        \qquad
        \tr(UZU^\dagger)=\tr(Z)
        \Longleftrightarrow U^\dagger UZ=Z.
    \]
    For every $W\in\mathbb C^{r\times r}$, one also has
    $\tr(U^\dagger WU)=\tr(W)$.
\end{lemma}

\begin{proof}\leanok
    Put $P=U^\dagger U$.  The matrices $P$ and $1-P$ are orthogonal
    projections, and
    \[
        \tr(UZU^\dagger)=\tr(PZ),
        \qquad
        \tr(Z)-\tr(PZ)
          =\tr((1-P)Z)\geq0.
    \]
    The last trace vanishes precisely when $(1-P)Z=0$.  The formula for $W$
    follows by cyclicity of trace and $UU^\dagger=1$.
\end{proof}

\begin{definition}[Matrices before active-sector compression]
    \label{def:mpdo_vertical_sector_precompression}
    \lean{MPOTensor.precompressionVerticalSectorT}
    \lean{MPOTensor.precompressionVerticalSectorS}
    \leanok
    Let $Z_T(X)$ be the matrix obtained by applying the normalized one-site
    embedding, expanding by $U_1^\dagger$, applying $T$, and relabelling the
    two-site indices, before compression by $U_2$.  Define $Z_S(Y)$ in the
    reverse direction, before compression by $U_1$.
\end{definition}

\begin{lemma}[Positivity and trace before active-sector compression]
    \label{lem:mpdo_vertical_sector_precompression_positive_trace}
    \lean{MPOTensor.precompressionVerticalSectorT_posSemidef}
    \lean{MPOTensor.precompressionVerticalSectorS_posSemidef}
    \lean{MPOTensor.trace_precompressionVerticalSectorT}
    \lean{MPOTensor.trace_precompressionVerticalSectorS}
    \lean{MPOTensor.verticalSectorTrace_transportedVerticalSectorT}
    \lean{MPOTensor.verticalSectorTrace_transportedVerticalSectorS}
    \leanok
    \uses{def:mpdo_transported_vertical_sector_maps,
      def:mpdo_positive_vertical_sector_family,
      def:mpdo_vertical_sector_precompression,
      lem:mpdo_normalized_vertical_sector_positive_trace}
    Suppose that the physical maps are completely positive and trace
    preserving and $U_iU_i^\dagger=1$.  For positive sector families
    $X\in\mathcal A_1$ and $Y\in\mathcal A_2$, the matrices $Z_T(X)$ and
    $Z_S(Y)$ are positive and
    \[
        \tr(Z_T(X))
          =\Trsec(X),
        \qquad
        \tr(Z_S(Y))
          =\Trsec(Y).
    \]
    After the final compression and sectorwise partial trace,
    \[
        \Trsec(\widetilde T(X))
          =\tr(U_2Z_T(X)U_2^\dagger),
        \qquad
        \Trsec(\widetilde S(Y))
          =\tr(U_1Z_S(Y)U_1^\dagger).
    \]
\end{lemma}

\begin{proof}\leanok
    Positivity follows successively from the normalized embedding, expansion
    by an adjoint, the physical map, relabelling, and partial trace.  The trace
    identities follow from the preceding two lemmas and trace preservation of
    the physical maps.
\end{proof}

% Local fix (zero-sector complement): the transported maps come from
% arXiv:1606.00608, Appendix C.4, lines 1955--1980, but the corrective trace-loss
% statement is recorded in docs/paper-gaps/cpgsv17_vertical_isometry_zero_sector.tex.
\begin{theorem}[Positivity and trace loss of the transported sector maps]
    \label{thm:mpdo_transported_vertical_sector_trace_loss}
    \lean{MPOTensor.transportedVerticalSectorT_posSemidef}
    \lean{MPOTensor.transportedVerticalSectorS_posSemidef}
    \lean{MPOTensor.transportedVerticalSectorT_trace_le}
    \lean{MPOTensor.transportedVerticalSectorS_trace_le}
    \lean{MPOTensor.transportedVerticalSectorT_trace_eq_iff}
    \lean{MPOTensor.transportedVerticalSectorS_trace_eq_iff}
    \leanok
    \uses{def:mpdo_transported_vertical_sector_maps,
      def:mpdo_positive_vertical_sector_family,
      def:mpdo_vertical_sector_precompression,
      lem:mpdo_normalized_retained_vertical_sector_embedding,
      lem:mpdo_retained_vertical_sector_partial_trace,
      lem:mpdo_coisometric_compression_trace,
      lem:mpdo_vertical_sector_precompression_positive_trace}
    Suppose that the multiplicity spaces are nonzero, their diagonal weights
    are positive, the physical maps $T$ and $S$ are completely positive and
    trace preserving, and the vertical coordinate maps satisfy
    $U_iU_i^\dagger=1$.  Then $\widetilde T$ and $\widetilde S$ preserve
    pointwise positive semidefiniteness.  For pointwise positive families
    $X\in\mathcal A_1$ and $Y\in\mathcal A_2$, their total sector traces obey
    \[
       \Trsec(\widetilde T(X))
          \leq \Trsec(X),
       \qquad
       \Trsec(\widetilde S(Y))
          \leq \Trsec(Y).
    \]
    The trace loss vanishes precisely when
    \[
       U_2^\dagger U_2 Z_T(X)=Z_T(X),
       \qquad
       U_1^\dagger U_1 Z_S(Y)=Z_S(Y),
    \]
    respectively.

    This result concerns positivity and trace on positive elements.  It does
    not assert a complete-positivity structure for the finite products, a
    trace-preserving extension to their full ambient matrix spaces, or an
    inverse relation between the two transported maps.
\end{theorem}

\begin{proof}\leanok
    The normalized diagonal embedding and the sectorwise partial trace
    preserve positivity.  Before the final compression, trace preservation of
    the physical map and $U_iU_i^\dagger=1$ give the original sector trace.
    For a positive matrix $Z$ and a coisometry $U$, put
    $P=U^\dagger U$.  Then
    \[
       \tr(UZU^\dagger)
        =\tr(PZ)
        \leq\tr(Z).
    \]
    Since $1-P$ is an orthogonal projection, equality is equivalent to
    $\tr((1-P)Z)=0$, hence to $(1-P)Z=0$.
\end{proof}
